% ROZDZIAŁ 24 — TWIERDZENIA POKRYCIOWE
\chapter{Twierdzenia pokryciowe Vitaliego i Besicovitcha}
\label{ch:covering-24}

\section{Motywacja i terminologia}
W poprzednich rozdziałach zbiory mierzalne przybliżaliśmy
zbiorami otwartymi i zwartymi. W zastosowaniach do
różniczkowania całek i miar potrzebujemy mocniejszego narzędzia:
z dowolnej rodziny małych kul należy wybrać takie kule, których
nakładanie można kontrolować. W tym rozdziale podajemy trzy
odmienne rezultaty: lemat $5r$, twierdzenie Vitaliego i twierdzenie
Besicovitcha. Nie wolno ich utożsamiać: pierwsze dotyczy
\emph{powiększonych} kul, drugie pokrycia \emph{z dokładnością do
zbioru miary zero}, a trzecie pokrycia kulami \emph{bez powiększania},
z ograniczoną krotnością.

\begin{definicja}[Kule i krotność pokrycia]
\label{def:balls-24}
Niech $(X,d)$ będzie przestrzenią metryczną. Piszemy
$B(x,r)=\{y\in X:d(x,y)<r\}$ dla $r>0$.
Jeśli $B=B(x,r)$ i $a>0$, to $aB:=B(x,ar)$.
Rodzina $(B_i)_{i\in I}$ jest \emph{parami rozłączna}, jeżeli
$B_i\cap B_j=\varnothing$ dla $i\ne j$.
Jej \emph{krotność w punkcie} $x$ to liczba indeksów $i$ takich,
że $x\in B_i$ (dopuszczamy $+\infty$).
Funkcja charakterystyczna dowolnego zbioru $A$ jest w całym
skrypcie oznaczana przez $\chi_A$.
\end{definicja}

\begin{uwaga}
Twierdzenia sformułowane dla kul w $\mathbb R^d$ odnoszą się
zawsze do metryki euklidesowej i miary Lebesgue'a $\lambda_d$.
Przypadki metryczne bez miary rozpatrujemy oddzielnie.
Symbol $v_d:=\lambda_d(B(0,1))$ oznacza objętość kuli jednostkowej;
$\lambda_d(B(x,r))=v_dr^d$.
\end{uwaga}

\section{Lemat $5r$: wybór kul bez zakładania miary}
\begin{twierdzenie}[Lemat o pięciokrotnym powiększeniu]
\label{thm:5r-24}
Niech $(X,d)$ będzie przestrzenią metryczną, a $\mathcal B$
niepustą rodziną otwartych kul o promieniach ograniczonych
z góry przez pewną liczbę $R<\infty$.
Istnieje parami rozłączna podrodzina $\mathcal D\subseteq\mathcal B$
taka, że
\[
 \bigcup_{B\in\mathcal B}B\subseteq\bigcup_{D\in\mathcal D}5D.
\]
Jeżeli $X$ jest ośrodkowa, rodzina $\mathcal D$ jest co najwyżej
przeliczalna.
\end{twierdzenie}
\begin{proof}
Rozdzielmy kule na klasy według promienia:
$\mathcal B_k=\{B(x,r)\in\mathcal B:
2^{-k-1}R<r\le2^{-k}R\}$, $k=0,1,\ldots$.
W klasie $\mathcal B_0$ wybieramy maksymalną względem inkluzji
rodzinę parami rozłącznych kul. Indukcyjnie, po dokonaniu wyborów
z klas wcześniejszych, w $\mathcal B_k$ wybieramy maksymalną
rodzinę parami rozłącznych kul, rozłącznych również z kulami
wybranymi wcześniej. Maksymalność można uzyskać z lematu
Zorna: uporządkowujemy przez inkluzję rodziny spełniające
warunek rozłączności, a suma łańcucha pozostaje taką rodziną.
Niech $\mathcal D$ będzie sumą wszystkich wybranych rodzin.

Weźmy kulę $B=B(x,r)\in\mathcal B_k$, która nie została
wybrana. Na odpowiednim etapie była albo niedopuszczalna
wskutek przecięcia z wcześniej wybraną kulą, albo przecinała
kulę z maksymalnego wyboru w swojej klasie. W obu przypadkach
istnieje $D=B(y,s)\in\mathcal D$ taka, że $B\cap D\ne\varnothing$
i $s>2^{-k-1}R\ge r/2$. Dla każdego $z\in B$ mamy
\[
 d(z,y)\le d(z,x)+d(x,y)<r+(r+s)=2r+s<5s.
\]
Stąd $B\subseteq5D$. Teza jest oczywista także dla kul
wybranych, bo $D\subseteq5D$.

Jeżeli $X$ ma przeliczalny zbiór gęsty $Q$, to każda wybrana
kula zawiera punkt $Q$. Kule z $\mathcal D$ są rozłączne,
więc przyporządkowanie im różnych punktów $Q$ dowodzi
przeliczalności $\mathcal D$.
\end{proof}

\begin{wniosek}[Oszacowanie miary sumy kul]
\label{cor:5r-volume-24}
Niech $\mathcal B$ będzie rodziną kul euklidesowych o promieniach
ograniczonych z góry. Istnieje parami rozłączna podrodzina
$(B_i)$ taka, że
\[
 \lambda_d\Bigl(\bigcup_{B\in\mathcal B}B\Bigr)
 \le 5^d\sum_i\lambda_d(B_i).
\]
\end{wniosek}
\begin{proof}
Stosujemy twierdzenie~\ref{thm:5r-24}, przeliczalną
podaddytywność oraz równość $\lambda_d(5B_i)=5^d\lambda_d(B_i)$.
Ponieważ suma rozłącznych kul jest mierzalna, szereg po prawej
jest dobrze określony również wtedy, gdy ma wartość $+\infty$.
\end{proof}

\begin{uwaga}[Inne stałe]
W przypadku kul w $\mathbb R^d$ często formułuje się wariant
$3r$, jeżeli promienie porządkuje się dokładniej. Stała $5$
jest wygodna, gdy nie zakłada się osiągania supremum promieni;
jej konkretna wartość nie jest istotna w większości zastosowań.
Nie należy jednak z lematu $5r$ wyciągać wniosku, że same
wybrane kule pokrywają całą początkową rodzinę.
\end{uwaga}

\begin{center}
\begin{tikzpicture}[scale=0.82]
  \draw[gray!80,thick] (0,0) circle (0.65);
  \draw[gray!80,thick] (1.18,0.18) circle (0.66);
  \draw[gray!80,thick] (3.1,0) circle (0.55);
  \fill (0,0) circle (1.6pt);
  \fill (1.18,0.18) circle (1.6pt);
  \fill (3.1,0) circle (1.6pt);
  \draw[blue!70!black,thick] (0,0) circle (0.65);
  \draw[blue!70!black,thick] (3.1,0) circle (0.55);
  \draw[blue!50!black,densely dashed] (0,0) circle (1.75);
  \draw[blue!50!black,densely dashed] (3.1,0) circle (1.45);
  \node[font=\small] at (1.5,-2.2) {Wybrane kule i ich powiększenia (schemat)};
  \node[font=\scriptsize,align=center] at (1.2,1.05) {kula\\odrzucona};
\end{tikzpicture}
\end{center}
\noindent\textit{Idea lematu $5r$:} kula odrzucona przecina jedną
z wybranych kul o porównywalnym promieniu, dlatego jest
zawarta w jej pięciokrotnym powiększeniu. Okręgi
kreskowane mają na rysunku promienie poglądowe, nie
w skali $5{:}1$.

\section{Pokrycia drobne i twierdzenie Vitaliego}
\begin{definicja}[Drobne pokrycie Vitaliego]
\label{def:vitali-fine-24}
Niech $E\subseteq\mathbb R^d$ i $\mathcal V$ będzie rodziną
otwartych kul euklidesowych. Mówimy, że $\mathcal V$ jest
\emph{drobne w każdym punkcie $E$}, jeżeli dla wszystkich
$x\in E$ i $\varepsilon>0$ istnieje kula
$B(x,r)\in\mathcal V$ ze środkiem w $x$ i $0<r<\varepsilon$.
Wymaganie środka w punkcie $x$ jest tutaj częścią definicji.
\end{definicja}

\begin{twierdzenie}[Pokryciowe twierdzenie Vitaliego]
\label{thm:vitali-cover-24}
Niech $E\subseteq\mathbb R^d$ będzie ograniczony,
$O\subseteq\mathbb R^d$ otwarty i ograniczony, $E\subseteq O$,
a $\mathcal V$ rodziną kul drobną w każdym punkcie $E$.
Wówczas można wybrać skończoną lub przeliczalną rodzinę
parami rozłącznych kul $(B_j)$ z $\mathcal V$, spełniającą
$B_j\subseteq O$ oraz
\[
 \lambda_d^*\Bigl(E\setminus\bigcup_jB_j\Bigr)=0,
 \qquad \sum_j\lambda_d(B_j)\le\lambda_d(O).
\]
Zapis $\lambda_d^*$ oznacza zewnętrzną miarę Lebesgue'a;
nie zakładamy z góry mierzalności zbioru $E$.
\end{twierdzenie}
\begin{proof}
Ograniczamy rodzinę do tych kul, których domknięcia mieszczą
się w $O$. Nadal jest ona drobna w każdym punkcie $E$:
przy ustalonym $x\in E$ wybieramy promień mniejszy od
odległości $x$ od $\mathbb R^d\setminus O$.
Wszystkie takie kule mieszczą się w $O$, więc ich promienie
są jednostajnie ograniczone.

Budujemy indukcyjnie rozłączne kule. Jeżeli po wybraniu
$B_1,\ldots,B_{n-1}$ pozostaje choć jedna kula naszej rodziny
rozłączna ze wszystkimi wcześniej wybranymi, oznaczmy przez
$s_n$ supremum promieni takich kul. Wybierzmy jedną z nich,
$B_n=B(x_n,r_n)$, o promieniu $r_n>s_n/2$. Jeżeli nie ma
dopuszczalnych kul, kończymy. Konstrukcja daje
$\sum_jv_dr_j^d\le\lambda_d(O)<\infty$, a zatem
$r_j\to0$ w przypadku nieskończonego ciągu.

Niech $N$ będzie dowolną dodatnią liczbą całkowitą.
Pokażemy, że poza zbiorem
$Z=\bigcup_j\partial B_j$ zachodzi inkluzja
\begin{equation}\label{eq:vitali-tail-24}
 E\setminus\bigcup_j\overline{B_j}
 \subseteq\bigcup_{j>N}5B_j.
\end{equation}
Rzeczywiście, ustalmy $x\in E\setminus\bigcup_j\overline{B_j}$.
Jest on w dodatniej odległości od skończonej sumy
$\bigcup_{j\le N}\overline{B_j}$. Z drobności pokrycia
wybieramy kulę $B=B(x,r)\in\mathcal V$ o domknięciu
w $O$, rozłączną z $B_1,\ldots,B_N$.
Jeżeli konstrukcja zakończyła się przed krokiem $N+1$,
taka kula przeczyłaby warunkowi zakończenia. Jeżeli zaś
konstrukcja trwa nieskończenie długo, $B$ musi przecinać
którąś z późniejszych kul. Gdyby bowiem była rozłączna
ze wszystkimi $B_j$, pozostawałaby dopuszczalna na
każdym etapie. Otrzymalibyśmy $s_j\ge r$ i
$r_j>s_j/2\ge r/2$ dla wszystkich $j>N$,
sprzecznie z $\sum_jr_j^d<\infty$.
Niech zatem $j>N$ będzie pierwszym indeksem, dla
którego $B\cap B_j\ne\varnothing$.
Przed krokiem $j$ kula $B$ pozostawała dopuszczalna,
zatem $r\le s_j<2r_j$. Odległość centrów jest mniejsza
od $r+r_j$, a więc $d(x,x_j)<3r_j<5r_j$.
Otrzymaliśmy $x\in5B_j$, co daje \eqref{eq:vitali-tail-24}.

Każda sfera $\partial B_j$ ma zerową miarę Lebesgue'a
(co wynika na przykład z zawierania jej w powłoce kulistej
o dowolnie małej objętości). Zatem $\lambda_d(Z)=0$.
Ze \eqref{eq:vitali-tail-24}, podaddytywności i skalowania
miary wynika
\[
 \lambda_d^*\Bigl(E\setminus\bigcup_jB_j\Bigr)
 \le\lambda_d(Z)+5^d\sum_{j>N}\lambda_d(B_j).
\]
Po przejściu $N\to\infty$ prawa strona dąży do zera.
Suma objętości nie przekracza $\lambda_d(O)$, bo kule są
parami rozłączne i zawarte w $O$. To kończy dowód.
\end{proof}

\begin{uwaga}[Dwa niezbędne ograniczenia w sformułowaniu]
Drobność jest warunkiem \emph{dla każdego punktu} $E$,
a nie tylko prawie wszędzie. Ograniczoność $O$ zapewnia
sumowalność objętości wybranych kul. Dla nieograniczonego
zbioru można podzielić przestrzeń na rozłączne warstwy
o skończonej mierze i zastosować argument lokalnie,
trzeba jednak zadbać o rozłączność wyborów pomiędzy warstwami.
Powyższa wersja wystarcza do różniczkowania miary Lebesgue'a.
\end{uwaga}

\begin{wniosek}[Przybliżanie zbioru rozłącznymi kulami]
\label{cor:vitali-open-24}
Jeżeli $O\subseteq\mathbb R^d$ jest niepustym,
ograniczonym zbiorem otwartym, to istnieją parami rozłączne
kule $B_j\subseteq O$ takie, że
$\lambda_d(O\setminus\bigcup_jB_j)=0$.
\end{wniosek}
\begin{proof}
Stosujemy twierdzenie~\ref{thm:vitali-cover-24} do $E=O$
i rodziny wszystkich kul o domknięciu zawartym w $O$.
Ta rodzina jest drobna w każdym punkcie $O$.
\end{proof}

\section{Geometria ograniczonej krotności: lemat kątowy}
\begin{lemat}[Skończoność konfiguracji Besicovitcha]
\label{lem:besicovitch-angle-24}
Niech $B_i=B(x_i,r_i)$, $i=1,\ldots,m$, będą kulami
w $\mathbb R^d$, których część wspólna zawiera punkt $p$.
Załóżmy ponadto, że dla $i\ne j$ zachodzi
$x_i\notin B_j$ i $x_j\notin B_i$.
Wówczas $m\le N_d$, gdzie $N_d$ zależy wyłącznie od $d$;
można przyjąć $N_d=3^d$.
\end{lemat}
\begin{proof}
Odrzucamy ewentualną kulę o środku $p$: skoro $p$ leży
we wszystkich kulach, taka kula wymusiłaby $m=1$.
Dla pozostałych stawiamy $u_i=(x_i-p)/|x_i-p|\in S^{d-1}$.
Ponieważ $p\in B_i$, mamy $|x_i-p|<r_i$.
Z warunku braku centrów w innych kulach wynika
$|x_i-x_j|\ge\max\{r_i,r_j\}>
\max\{|x_i-p|,|x_j-p|\}$.
Z twierdzenia cosinusów
\[
 2\langle x_i-p,x_j-p\rangle
 =|x_i-p|^2+|x_j-p|^2-|x_i-x_j|^2
 <\min\{|x_i-p|^2,|x_j-p|^2\}.
\]
W każdym przypadku z tej
nierówności otrzymujemy
$\langle u_i,u_j\rangle<1/2$, a zatem
$|u_i-u_j|>1$. Kulki $B(u_i,1/2)$ w $\mathbb R^d$
są parami rozłączne i leżą w $B(0,3/2)$.
Porównanie objętości daje
$m v_d(1/2)^d\le v_d(3/2)^d$, skąd $m\le3^d$.
\end{proof}

\begin{twierdzenie}[Skończona postać lematu Besicovitcha]
\label{thm:besicovitch-finite-24}
Niech $x_1,\ldots,x_M$ będą punktami $\mathbb R^d$,
a $r_1,\ldots,r_M>0$. Istnieje podrodzina kul
$B(x_i,r_i)$, która pokrywa wszystkie punkty
$x_1,\ldots,x_M$ i której krotność w żadnym punkcie
$\mathbb R^d$ nie przekracza $N_d=3^d$.
\end{twierdzenie}
\begin{proof}
Porządkujemy kule nierosnąco według promienia (dla remisów
wybieramy dowolny porządek). Przechodzimy po kolei przez
ich listę i wybieramy kulę, jeśli jej środek nie należy
do sumy kul wybranych wcześniej. Każdy środek jest ostatecznie
pokryty: albo przez własną wybraną kulę, albo przez jedną
z poprzednich. Jeśli wybrano $B_i$ wcześniej niż $B_j$,
to $x_j\notin B_i$, a stąd $|x_i-x_j|\ge r_i\ge r_j$
i także $x_i\notin B_j$. Wybrane kule spełniają więc
warunek lematu~\ref{lem:besicovitch-angle-24}.
Dla dowolnego punktu $p$ kule wybrane zawierające $p$
tworzą konfigurację z tego lematu, zatem jest ich najwyżej
$N_d$.
\end{proof}

\begin{wniosek}[Oszacowanie sumy miar w rodzinie skończonej]
\label{cor:finite-besicovitch-measure-24}
W założeniach twierdzenia~\ref{thm:besicovitch-finite-24}
niech $B_{i_1},\ldots,B_{i_s}$ będzie wybraną podrodziną.
Dla dowolnej dodatniej miary borelowskiej $\rho$
na $\mathbb R^d$ zachodzi
\[
 \sum_{j=1}^s\rho(B_{i_j})
 \le N_d\,\rho\Bigl(\bigcup_{j=1}^s B_{i_j}\Bigr).
\]
Teza pozostaje prawdziwa, gdy prawa strona jest nieskończona.
\end{wniosek}
\begin{proof}
Definicja krotności i skończona addytywność całki dają
\[
 \sum_{j=1}^s\rho(B_{i_j})
 =\int_{\mathbb R^d}\sum_{j=1}^s\chi_{B_{i_j}}(x)
 \,\mathrm d\rho(x)
 \le N_d\int_{\mathbb R^d}
 \chi_{\cup_jB_{i_j}}(x)\,\mathrm d\rho(x).
\]
To ostatnie wyrażenie jest równe prawej stronie tezy.
\end{proof}

\begin{twierdzenie}[Twierdzenie pokryciowe Besicovitcha -- pełny dowód]
\label{thm:besicovitch-24}
Niech $d\ge1$, niech $E\subseteq\mathbb R^d$ będzie
zbiorem ograniczonym i niech każdemu $x\in E$
przyporządkowano promień $r_x\in(0,R]$, gdzie
$R<\infty$. Wówczas z rodziny kul $\{B(x,r_x):x\in E\}$
można wybrać co najwyżej przeliczalną podrodzinę
$(B_j)_{j\in J}$ taką, że
\begin{equation}\label{eq:besicovitch-infinite-24}
 E\subseteq\bigcup_{j\in J}B_j,\qquad
 \sum_{j\in J}\chi_{B_j}(z)\le15^d
 \quad(z\in\mathbb R^d).
\end{equation}
Jeżeli dla każdego środka $x$ zadano kilka dopuszczalnych
kul o promieniach nie większych niż $R$, wolno najpierw
wybrać jedną z nich i zastosować twierdzenie.
Nie zakładamy mierzalności zbioru $E$.
\end{twierdzenie}
\begin{proof}
\emph{Krok 1: podział na skale.}
Dla $k=0,1,\ldots$ oznaczmy $a_k=2^{-k-1}R$ i
$E_k^0=\{x\in E:a_k<r_x\le2a_k\}$. Zbiory
$E_k^0$ tworzą rozłączny podział $E$.
Będziemy kolejno wybierać skończone zbiory środków
$S_k\subseteq E_k^0$; wybranej rodzinie w skali $k$
odpowiadają kule $\{B(x,r_x):x\in S_k\}$.

Przypuśćmy, że dokonano już wyboru dla skal
$0,\ldots,k-1$, i niech $O_{k-1}$ oznacza sumę
wszystkich wcześniej wybranych kul (dla $k=0$
przyjmujemy $O_{-1}=\varnothing$).
Niech $E_k=E_k^0\setminus O_{k-1}$. Spośród środków
$E_k$ wybieramy maksymalny zbiór $S_k$ o własności
\begin{equation}\label{eq:besicovitch-separated-24}
 |x-y|\ge a_k\quad(x,y\in S_k,\ x\ne y).
\end{equation}
Taki zbiór można zbudować zachłannie, bez lematu
Zorna: zaczynamy od dowolnego środka, a następnie
dodajemy kolejny, jeśli jego odległość od każdego
wcześniej wybranego jest co najmniej $a_k$.
Ponieważ $E$ jest ograniczony, każde dwa punkty
wybranego zbioru są oddalone o co najmniej $a_k$,
a kule $B(x,a_k/2)$ o tych środkach są rozłączne
i leżą we wspólnej ograniczonej kuli. Porównanie
ich objętości ogranicza z góry liczbę kroków;
procedura kończy się więc po skończenie wielu
wyborach, jeżeli $E_k\ne\varnothing$.

Z maksymalności wynika, że każdy $x\in E_k$
albo należy do $S_k$, albo istnieje $y\in S_k$
z $|x-y|<a_k<r_y$. Toteż wszystkie środki z
$E_k$ należą do kul wybranych w skali $k$.
W konsekwencji każdy $x\in E$ zostaje pokryty:
jeżeli $x\in E_k^0$, to albo znajdował się już
w $O_{k-1}$, albo zostaje pokryty w kroku $k$.
Suma $\bigcup_{k\ge0}S_k$ jest przeliczalna,
gdyż wszystkie $S_k$ są skończone.

\emph{Krok 2: liczba kul jednej skali.}
Ustalmy $z\in\mathbb R^d$ i skalę $k$.
Jeżeli $B(x,r_x)$ jest wybraną kulą skali $k$
zawierającą $z$, to $|x-z|<r_x\le2a_k$.
Z \eqref{eq:besicovitch-separated-24} wynika, że
kule $B(x,a_k/2)$ o takich środkach są rozłączne,
a każda z nich mieści się w $B(z,5a_k/2)$.
Porównując ich objętości, otrzymujemy oszacowanie
\begin{equation}\label{eq:besicovitch-single-scale-24}
 \#\{x\in S_k:z\in B(x,r_x)\}\le5^d.
\end{equation}

\emph{Krok 3: liczba skal występujących w punkcie.}
Spośród każdej skali $k$, dla której jakaś wybrana
kula zawiera $z$, wybierzmy jedną taką kulę
$B(x_k,r_{x_k})$. Gdy $k<\ell$, środek
$x_\ell$ nie należy do żadnej kuli wybranej w
skali $k$ (tak zdefiniowano $E_\ell$), więc
$|x_k-x_\ell|\ge r_{x_k}$. Ponadto
$r_{x_k}>a_k\ge r_{x_\ell}$, a zatem
$x_k\notin B(x_\ell,r_{x_\ell})$.
Tak wybrane kule wszystkich występujących skal
mają wspólny punkt $z$ i spełniają oba warunki
wyłączenia środków z
Lematu~\ref{lem:besicovitch-angle-24}.
Dowolna \emph{skończona} liczba tych skal jest
więc ograniczona przez $3^d$. Gdyby skal było
więcej niż $3^d$, otrzymalibyśmy sprzeczność
już dla pewnej skończonej podrodziny.

Każda skala wnosi najwyżej $5^d$ kul
zawierających $z$, a skal jest najwyżej $3^d$.
Daje to drugą nierówność w
\eqref{eq:besicovitch-infinite-24} i kończy dowód.
\end{proof}

\begin{uwaga}[Dlaczego metoda działa dla rodzin nieskończonych?]
Kluczowe jest oddzielenie dwóch rodzajów nakładania.
W jednej skali wystarcza pakowanie oddalonych
środków, natomiast pomiędzy skalami obowiązuje
silniejsza reguła: środek późniejszej kuli nie leży
w żadnej kuli wcześniej wybranej. Dopiero
połączenie obu argumentów daje jednostajną
krotność $15^d$. Nie przechodzimy tu od pokryć
skończonych do granicy i nie powiększamy żadnej kuli.
Stała nie jest optymalna; istotna jest jej
zależność wyłącznie od wymiaru.
\end{uwaga}

\begin{wniosek}[Besicovitch dla dodatnich miar borelowskich]
\label{cor:besicovitch-measure-infinite-24}
Przy założeniach Twierdzenia~\ref{thm:besicovitch-24},
dla dowolnej dodatniej miary borelowskiej $\rho$
na $\mathbb R^d$ wybrana rodzina kul spełnia
\[
 \sum_{j\in J}\rho(B_j)
 \le15^d\,\rho\Bigl(\bigcup_{j\in J}B_j\Bigr).
\]
\end{wniosek}
\begin{proof}
Stosujemy twierdzenie Tonellego do
$\sum_j\chi_{B_j}$ i korzystamy z oszacowania
krotności w \eqref{eq:besicovitch-infinite-24}.
\end{proof}

\section{Zastosowania do oszacowań typu maksymalnego}
\begin{definicja}[Funkcja maksymalna Hardy'ego--Littlewooda]
\label{def:HL-max-24}
Dla $f\in L^1_{\mathrm{loc}}(\mathbb R^d)$ definiujemy
\emph{centrowaną funkcję maksymalną}
\[
 Mf(x):=\sup_{r>0}\frac1{\lambda_d(B(x,r))}
 \int_{B(x,r)}|f(y)|\,\mathrm d\lambda_d(y)
 \in[0,+\infty].
\]
Funkcja $Mf$ jest mierzalna, co wynika z faktu, że dla
ustalonego $r$ średnia po $B(x,r)$ jest funkcją ciągłą
zmiennej $x$, oraz z możliwości ograniczenia supremum
do promieni wymiernych dodatnich. Ciągłość średniej wynika
z ciągłości translacji w $L^1$ na odpowiednich zbiorach
zwartych, a ciągłość względem $r>0$ z bezwzględnej
ciągłości całki po kulach.
\end{definicja}

\begin{twierdzenie}[Nierówność słabego typu $(1,1)$]
\label{thm:HL-weak-24}
Jeśli $f\in L^1(\mathbb R^d)$, to dla każdego $t>0$
\[
 \lambda_d\bigl(\{x\in\mathbb R^d:Mf(x)>t\}\bigr)
 \le\frac{5^d}{t}\int_{\mathbb R^d}|f(x)|\,
 \mathrm d\lambda_d(x).
\]
\end{twierdzenie}
\begin{proof}
Ustalmy zwarty zbiór $K\subseteq\{Mf>t\}$.
Dla każdego $x\in K$ wybieramy promień $r_x>0$ taki,
że $\int_{B(x,r_x)}|f|\,\mathrm d\lambda_d>
 t\lambda_d(B(x,r_x))$. Promienie takich kul są
ograniczone z góry: ponieważ ich objętość jest mniejsza
od $\|f\|_1/t$, zachodzi $r_x^d<\|f\|_1/(tv_d)$.
Z lematu $5r$ wybieramy rozłączne kule $(B_j)$, których
pięciokrotne powiększenia pokrywają $K$. Wówczas
\[
 \lambda_d(K)\le5^d\sum_j\lambda_d(B_j)
 <\frac{5^d}{t}\sum_j\int_{B_j}|f|\,\mathrm d\lambda_d
 \le\frac{5^d}{t}\|f\|_1.
\]
Zbiór $\{Mf>t\}$ jest otwarty: dla dowolnego $x$
warunek średniej większej od $t$ zachowuje się przy małej
zmianie środka kuli dzięki ciągłości całki z $|f|$.
Jego miara jest supremum miar zwartych podzbiorów,
na mocy regularności miary Lebesgue'a. Biorąc supremum
po $K$, kończymy dowód.
\end{proof}

\begin{wniosek}[Kontrola zbioru, na którym średnie są duże]
\label{cor:HL-level-24}
Dla $f\in L^1(\mathbb R^d)$ oraz $a>0$ funkcja $Mf$
jest skończona prawie wszędzie, a
$\lambda_d(\{Mf>a\})\le5^d\|f\|_1/a$.
\end{wniosek}
\begin{proof}
Oszacowanie wynika z twierdzenia~\ref{thm:HL-weak-24}.
Przechodząc do granicy $a\to\infty$, dostajemy
$\lambda_d(\{Mf=+\infty\})=0$.
\end{proof}

\begin{uwaga}[Co zostawiamy do następnego rozdziału?]
W rozdziale~\ref{ch:covering-24} skonstruowaliśmy narzędzie
pokryciowe i udowodniliśmy nierówność słabego typu.
W kolejnym rozdziale użyjemy ich do dowodu twierdzenia
Lebesgue'a o różniczkowaniu całek i miar, bez zakładania
go w argumentach niniejszego rozdziału.
\end{uwaga}
